% !TEX root = ../main.tex
\chapter{Language Models}
\label{ch:language_models}

Language models are the Transformer's \cite{vaswani_attention_2017}
historical and contemporary main use. \Lucid\ provides the
language-model side of the zoo: BERT and its derivatives
(encoder-only) \cite{devlin_bert_2019, liu_roberta_2019,
lan_albert_2020}, GPT and its derivatives (decoder-only)
\cite{radford_gpt2_2019, brown_gpt3_2020, touvron_llama_2023},
T5 (encoder--decoder) \cite{raffel_t5_2020}, RoFormer with
rotary positional embeddings \cite{su_roformer_2021}, and
small LoRA \cite{hu_lora_2021} adapters for efficient
fine-tuning. The chapter describes the architectural variants,
the tokenisation pipeline, the inference-time KV-cache
optimisation, and Apple-silicon-specific implementation choices.

\section{The Transformer Recap}
\label{sec:language_models:transformer_recap}

A Transformer block, the universal building block, performs the
following on an input $X \in \mathbb{R}^{B \times T \times D}$:

\begin{align}
  \tilde{X} &= X + \mathrm{MHA}(\mathrm{LN}(X)), \\
  Y &= \tilde{X} + \mathrm{FFN}(\mathrm{LN}(\tilde{X})),
\end{align}
where MHA is multi-head attention
(\chref{ch:attention}~Section~\ref{sec:attention:multi_head}) and FFN
is the position-wise feed-forward network:
\[
  \mathrm{FFN}(x) \;=\; W_2 \cdot \phi(W_1 x + b_1) + b_2,
\]
with $W_1 \in \mathbb{R}^{D_{\text{ff}} \times D}$ and
$W_2 \in \mathbb{R}^{D \times D_{\text{ff}}}$ and $\phi$ an
activation (GELU \cite{hendrycks_gelu_2016}, ReLU, or
SwiGLU \cite{shazeer_glu_2020}). Typically
$D_{\text{ff}} = 4D$.

The arrangement of LayerNorm inside the residual (called
\emph{pre-norm}, the modern preference) gives more stable
training than the original \emph{post-norm}. \Lucid\ defaults to
pre-norm for all language models.

\subsection{Position encoding}
\label{ssec:language_models:pos_encoding}

The Transformer is permutation-equivariant; positional
information must be injected. Three options:

\begin{itemize}
  \item \textbf{Sinusoidal} (original Transformer
    \cite{vaswani_attention_2017}): fixed sin/cos functions of
    position added to the input embedding.
  \item \textbf{Learned absolute} (BERT, GPT): a learnable
    $T_{\max} \times D$ matrix indexed by position. Limits to
    the training context length.
  \item \textbf{Rotary} (RoPE \cite{su_roformer_2021}): rotates
    the query and key vectors in 2D subspaces by an angle
    proportional to position. Naturally relative, extends past
    the training length.
\end{itemize}

\Lucid\ implements all three; the model factory chooses
appropriately (e.g., RoFormer uses RoPE, BERT uses learned
absolute).

\section{BERT (Encoder-only)}
\label{sec:language_models:bert}

BERT \cite{devlin_bert_2019} is a stack of Transformer encoder
blocks pretrained with masked language modelling. The encoder
attends bidirectionally (no causal mask); the [CLS] token's
final embedding is used for sentence-level tasks (classification,
NLI), and the per-token embeddings are used for token-level
tasks (NER, QA span extraction).

\subsection{Pretraining objective}
\label{ssec:language_models:bert_pretrain}

The masked-language-modelling loss masks 15\% of input tokens
and predicts them from the bidirectional context:
\[
  \mathcal{L}_{\text{MLM}} \;=\; -\sum_{i \in \mathcal{M}} \log p(x_i \mid x_{\setminus \mathcal{M}}),
\]
where $\mathcal{M}$ is the set of masked positions.
The next-sentence-prediction loss (original BERT) was found
unnecessary in later work \cite{liu_roberta_2019}; \Lucid's BERT
implementation supports MLM-only training.

\subsection{Variants}
\label{ssec:language_models:bert_variants}

The paper lists two sizes:
\begin{itemize}
  \item \code{bert\_base}: $L = 12$, $D = 768$, $H = 12$,
    110M params.
  \item \code{bert\_large}: $L = 24$, $D = 1024$, $H = 16$,
    340M params.
\end{itemize}

Additional variants \Lucid\ ships:
\begin{itemize}
  \item \code{roberta\_base}, \code{roberta\_large}:
    re-trained BERT with more data and longer training
    \cite{liu_roberta_2019}.
  \item \code{albert\_base}, \code{albert\_large},
    \code{albert\_xlarge}, \code{albert\_xxlarge}:
    parameter-shared BERT \cite{lan_albert_2020}.
\end{itemize}

\section{GPT (Decoder-only)}
\label{sec:language_models:gpt}

GPT \cite{radford_gpt2_2019, brown_gpt3_2020} is a stack of
Transformer decoder blocks pretrained with autoregressive
language modelling. The decoder attends causally: each token
attends only to itself and earlier tokens. The output at each
position is the predicted distribution over the next token.

\subsection{The causal mask}
\label{ssec:language_models:causal_mask}

Implemented as an additive mask in the attention's softmax:
\[
  M_{ij} \;=\;
  \begin{cases}
    0, & j \leq i, \\
    -\infty, & j > i.
  \end{cases}
\]
The mask is added to $QK^\top/\sqrt{d_h}$ before softmax;
the $-\infty$ entries zero their softmax probabilities.

On Apple silicon, the SDPA fused kernel
(\chref{ch:attention}~Section~\ref{sec:attention:sdpa}) supports
causal mode without materialising the mask:
\code{F.scaled\_dot\_product\_attention(q, k, v, is\_causal=True)}.
The kernel uses a constant-stride bit-pattern to apply the
mask, saving $O(T^2)$ memory.

\subsection{Pretraining objective}
\label{ssec:language_models:gpt_pretrain}

\[
  \mathcal{L}_{\text{LM}} \;=\; -\sum_{t=1}^T \log p(x_t \mid x_{<t}).
\]
This is straightforward cross-entropy with the model's
distribution at each position against the ground-truth next
token.

\subsection{Variants}
\label{ssec:language_models:gpt_variants}

\code{gpt2\_small} ($L = 12$, $D = 768$, $H = 12$, 124M params),
\code{gpt2\_medium} ($L = 24$, $D = 1024$, 355M params),
\code{gpt2\_large} ($L = 36$, $D = 1280$, 774M params),
\code{gpt2\_xlarge} ($L = 48$, $D = 1600$, 1.5B params)
\cite{radford_gpt2_2019}.

\Lucid\ also provides \code{llama\_7b} (and the
larger \code{llama\_13b}, \code{llama\_30b}, \code{llama\_65b}
configs that work but are too large to train on a single Mac).
LLaMA \cite{touvron_llama_2023} introduces:
\begin{itemize}
  \item RoPE positional encoding (no learnable position).
  \item RMSNorm \cite{zhang_rmsnorm_2019} instead of LayerNorm:
    $\mathrm{RMSNorm}(x) = x \cdot \mathrm{rsqrt}(\mathrm{mean}(x^2) + \epsilon) \cdot \gamma$.
  \item SwiGLU activation \cite{shazeer_glu_2020}:
    $\mathrm{SwiGLU}(x) = (x W_1) \odot \mathrm{SiLU}(x W_3) \cdot W_2$.
\end{itemize}

\section{T5 (Encoder--Decoder)}
\label{sec:language_models:t5}

T5 \cite{raffel_t5_2020} retains the original Transformer's
encoder--decoder architecture and reformulates every NLP task
as text-to-text:
\begin{itemize}
  \item Translation: ``translate English to German: That is
    good.'' $\to$ ``Das ist gut.''
  \item Classification: ``cola sentence: He is is.'' $\to$
    ``unacceptable''.
  \item Summarisation: ``summarize: \ldots'' $\to$ ``\ldots''
\end{itemize}

The encoder is bidirectional; the decoder attends causally to
itself and full to the encoder output (cross-attention). The
T5 architecture has minor differences from the vanilla
Transformer:
\begin{itemize}
  \item Pre-norm (RMSNorm placed before each block).
  \item Relative position bias added inside attention, not at
    the input.
  \item ReLU (not GELU) in the FFN.
  \item No bias terms in the linear layers.
\end{itemize}

Variants: \code{t5\_small} (60M), \code{t5\_base} (220M),
\code{t5\_large} (770M), \code{t5\_3b}, \code{t5\_11b}
(the larger sizes are inference-only on M-series Macs).

\section{Tokenisation}
\label{sec:language_models:tokenisation}

Language models do not consume raw bytes; the input is first
\emph{tokenised} into integer IDs. \Lucid\ provides:

\begin{itemize}
  \item \textbf{BPE} (Byte-Pair Encoding) \cite{sennrich_bpe_2016}:
    used by GPT-2, RoBERTa.
  \item \textbf{WordPiece}: used by BERT (similar to BPE but
    splits at character boundaries with \texttt{\#\#} markers).
  \item \textbf{SentencePiece} \cite{kudo_sentencepiece_2018}:
    used by T5, LLaMA. Whitespace-aware, language-agnostic.
\end{itemize}

\subsection{Vocabulary size}
\label{ssec:language_models:vocab_size}

Typical vocab sizes: 30K (BERT), 50K (GPT-2), 32K (LLaMA).
The vocab choice trades:
\begin{itemize}
  \item Sequence length: larger vocab $\to$ shorter sequences
    for the same text $\to$ less attention compute.
  \item Embedding matrix size: larger vocab $\to$ more params
    in the $V \times D$ embedding table.
\end{itemize}

\subsection{The tokeniser API}
\label{ssec:language_models:tokeniser_api}

\begin{lstlisting}[style=lucid,language=LucidPython]
from lucid.models.tokenizers import AutoTokenizer

tok = AutoTokenizer.from_pretrained("gpt2")
ids = tok.encode("hello world")  # list[int]
text = tok.decode(ids)           # str
tensor_in = tok.encode_batch(["hello", "world"], padding=True)
\end{lstlisting}

The tokenisers are implemented in pure Python with
no external dependencies (per Rule H4 in the project
documentation).

\section{Inference-Time KV Cache}
\label{sec:language_models:kv_cache}

Autoregressive decoding produces one token at a time. Naive
implementation recomputes all of attention's keys and values at
every step, which is $O(T^2)$ in the sequence length $T$.

The KV cache stores keys and values from previous steps; the new
token's query attends to the cached keys and values plus its
own. This reduces per-step compute from $O(T^2)$ to $O(T)$.

\subsection{Memory cost}
\label{ssec:language_models:kv_cache_memory}

For a model with $L$ layers, $H$ heads, head dimension $d_h$,
and current sequence length $T$:
\[
  \text{KV memory} \;=\; 2 \cdot L \cdot H \cdot d_h \cdot T \cdot \text{dtype\_size}.
\]
For LLaMA-7B ($L = 32$, $H d_h = 4096$), at $T = 2048$ tokens
in FP16: $2 \cdot 32 \cdot 4096 \cdot 2048 \cdot 2 = 1.07\,\mathrm{GB}$.
This memory cost is one of the main bottlenecks for long-context
inference.

\subsection{The Lucid KV cache API}
\label{ssec:language_models:kv_cache_api}

\begin{lstlisting}[style=lucid,language=LucidPython]
model = lucid.models.gpt2_small().eval()
prompt = tokenizer.encode("Hello, ")
input_ids = lucid.tensor([prompt])

# Initialise: returns logits and the empty cache, prefilled with prompt
logits, kv_cache = model.generate_step(input_ids, kv_cache=None)
next_token = lucid.argmax(logits[:, -1])

# Step: one token at a time, using and updating the cache
for _ in range(50):
    logits, kv_cache = model.generate_step(next_token, kv_cache=kv_cache)
    next_token = lucid.argmax(logits[:, -1])
    output.append(next_token.item())
\end{lstlisting}

The cache is a list of (key, value) tuples, one per layer. On
Apple silicon, the cache is allocated in unified memory (zero
copy from the CPU's perspective), so cache updates are cheap.

\section{LoRA: Low-Rank Adaptation}
\label{sec:language_models:lora}

LoRA \cite{hu_lora_2021} fine-tunes large language models by
adding small low-rank residual adapters to the existing
weight matrices, freezing the base weights. For a weight matrix
$W \in \mathbb{R}^{D_{\text{out}} \times D_{\text{in}}}$, the
LoRA-adapted forward is:
\[
  y \;=\; W x + B A x,
\]
where $A \in \mathbb{R}^{r \times D_{\text{in}}}$ and
$B \in \mathbb{R}^{D_{\text{out}} \times r}$ with $r \ll D$.
Only $A$ and $B$ are trained; $W$ is frozen. The parameter cost
of the adapter is $r(D_{\text{in}} + D_{\text{out}})$, vs
$D_{\text{in}} \cdot D_{\text{out}}$ for full fine-tuning ---
typically a 100--1000$\times$ reduction.

\subsection{Implementation}
\label{ssec:language_models:lora_impl}

\Lucid's \code{lucid.models.lora} module provides a wrapper:

\begin{lstlisting}[style=lucid,language=LucidPython]
from lucid.models.lora import inject_lora

model = lucid.models.gpt2_small()
inject_lora(
    model,
    targets=["c_attn", "c_proj"],   # Layers to adapt
    rank=8,
    alpha=16,
)
# Now only the LoRA params are trainable
for name, p in model.named_parameters():
    if "lora_" not in name:
        p.requires_grad = False
\end{lstlisting}

The adapter weights $A$ are initialised to zero and $B$ to a
small random init; the product $BA$ is initially zero, so
fine-tuning starts from the pretrained model.

\section{Performance on Apple Silicon}
\label{sec:language_models:perf}

Per-step training time on M3~Max, batch 16, sequence 512,
FP32 (no compile yet):

\begin{table}[t]
\centering
\small
\begin{tabular}{lrr}
\toprule
Model & Params (M) & Step time (ms) \\
\midrule
\code{bert\_base}        & 110   & 140 \\
\code{bert\_large}       & 340   & 410 \\
\code{roberta\_base}     & 125   & 145 \\
\code{gpt2\_small}       & 124   & 155 \\
\code{gpt2\_medium}      & 355   & 430 \\
\code{t5\_small}         & 60    & 110 \\
\code{t5\_base}          & 220   & 290 \\
\code{llama\_7b} (FP16, frozen) & 6740  & 1850 (inference) \\
\bottomrule
\end{tabular}
\caption[Language-model training throughput.]{Per-step training
time on M3~Max for selected language models at batch 16, sequence
512, FP32 (except LLaMA-7B, FP16 inference-only). LLaMA-7B is too
large to train on a single Mac; it is included for reference.
Encoder-only (BERT) models are slightly faster than decoder-only
(GPT-2) of similar size because the causal mask adds attention
overhead.}
\label{tab:language_models:perf}
\end{table}

\section{Inference Throughput}
\label{sec:language_models:inference_throughput}

Generation throughput (tokens per second) on M3~Max with KV cache:

\begin{table}[t]
\centering
\small
\begin{tabular}{lrr}
\toprule
Model & Dtype & Tokens/s \\
\midrule
\code{gpt2\_small}       & FP32 & 220 \\
\code{gpt2\_small}       & FP16 & 380 \\
\code{gpt2\_medium}      & FP16 & 140 \\
\code{llama\_7b}         & FP16 & 22  \\
\code{llama\_7b}         & INT4 & 38  \\
\bottomrule
\end{tabular}
\caption[Generation throughput.]{Token-per-second generation
throughput on M3~Max with KV cache. Single-batch decoding;
prefill of 256 tokens excluded from the measurement.}
\label{tab:language_models:inference_throughput}
\end{table}

\section{Compile-Time Wins}
\label{sec:language_models:compile_wins}

Language models benefit from \code{lucid.compile} most
substantially when training:
\begin{itemize}
  \item \code{gpt2\_small} (training): 155 ms eager $\to$ 95 ms
    compiled (1.63$\times$).
  \item \code{bert\_base}: 140 ms $\to$ 88 ms (1.59$\times$).
  \item \code{t5\_base}: 290 ms $\to$ 178 ms (1.63$\times$).
\end{itemize}
The wins come from the same source as ViT
(\chref{ch:vision_transformers}~Section~\ref{sec:vision_transformers:compile_wins}):
fused SDPA + many normalisations + low per-block overhead.

Inference benefits as well, but less dramatically (10--20\%)
because the KV-cache update is already a small op count.

\section{Pretrained Weights}
\label{sec:language_models:pretrained}

\Lucid\ hosts the standard pretrained checkpoints for BERT,
RoBERTa, GPT-2, T5, RoFormer. LLaMA weights are not redistributed
(licensing); the user obtains them from the official source and
loads via:

\begin{lstlisting}[style=lucid,language=LucidPython]
model = lucid.models.AutoModel.from_pretrained("bert_base")
# Standard hub fetch

model = lucid.models.llama.from_local("/path/to/llama_7b_weights/")
# User-supplied
\end{lstlisting}

\section{Summary and Forward References}
\label{sec:language_models:summary}

\Lucid's language-model zoo: BERT (encoder), GPT-2/LLaMA
(decoder), T5 (encoder--decoder), RoFormer (RoPE), LoRA
(parameter-efficient fine-tuning). All variants follow the
standard zoo family contract; KV caching is supported for
decoder-style inference. The Apple-silicon performance is
competitive for sub-1B-parameter models; LLaMA-7B inference
runs but training is infeasible on a single Mac.

The next chapter, \chref{ch:generative}, treats generative
models (GAN, VAE, diffusion). \chref{ch:pretrained} treats the
weights hub.

\begin{lucidnote}
For the user: \code{gpt2\_small} is the typical entry point for
language-model experiments; for fine-tuning, prefer LoRA
adapters (10$\times$ less GPU memory, comparable accuracy). For
inference, enable the KV cache via \code{model.generate\_step}.
\end{lucidnote}
